% Transcription — fonds Alexandre Grothendieck, shelfmark 10, batch 8 (pages 141-160).
% Established from the facsimile archives/batches/10/batch-08.pdf (a mirror of
% https://grothendieck.umontpellier.fr/10.pdf), per the skill
% transcribe-grothendieck.
%
% Pass: Opus 5 (claude-opus-5), 2026-09-12 — first pass, unchecked against the
% pages by a human.
%
% The batch opens the run the brown folder cover of page 140 announced by its
% second line, « Anneaux des représentations de certains groupes algébriques ».
% Page 141 is blank. Pages 142-148 are his own hand, fast, numbered by him
% (1) to (34); page 149, in pencil, is a four-item recapitulation of the same
% formalism.
% Pages 150-151 open a second run, « A(T) = Z(T^) ». Pages 152-154 are a third,
% on K(G), CK(G) and the characteristic homomorphism, breaking off at (5 bis).
% Page 155 is blank. Pages 156-160 are a fourth and last, on an abelian class C
% of k-modules and the complete linearly topologised algebra U it reconstructs;
% it is the most heavily drafted of the lot and its last four leaves carry long
% cancellations.
%
%   142  the categories A^G of finite-dimensional k-vector spaces with a
%        G-action; the contravariant A^u attached to u : G → G'; the exterior
%        tensor product M ⊠ M' = A^{p_1}(M) ⊗ A^{p_2}(M'); the internal ⊗
%        recovered through the diagonal Δ_G; induction announced.
%   143  I^u and its transitivity; A^u I^u(M) = M^{[G':G]}; the two functorial
%        homomorphisms; the injections u_{p_1…p_h} of products of symmetric
%        groups; M * M' = I^{u_{p,q}}(M ⊠ M') with its commutativity through an
%        α satisfying α u_{p,q} α^{-1} = u_{q,p} s_{p,q}, and its associativity.
%   144  T^n(M) = ⊗^n M seen in (A^G)^{S_n}, a non-additive functor, with the
%        boxed T^n(M + M') = Σ_{p+q=n} T^p(M) * T^q(M'); the compatibility
%        square with A^u; u_{p,q} : S_p × S_q → S_{pq} written out through the
%        α_{p,q}(σ), and T^m(T^n(M)) ≃ A^{u_{m,n}}(T^{mn}(M)).
%   145  one line: T^n(M ⊗ M') ≃ T^n(M) ⊗ T^n(M').
%   146  * is bi-additive with unit 1 = K in A^{S_0}; A(G) = K(A^G), a Chern
%        ring; A(u) contravariant, α ⊡ β, the induction I(u) additive but not
%        multiplicative, and (26) on the two composites.
%   147  A(σ_g) = I(σ_g) = id; the graded ring A^• = ∐ A^n with A^0 = Z;
%        A(u_{p,q})(α * β) = ((p+q)!/p!q!) α ⊡ β; T as a homomorphism into
%        1 + completion of A(G)^{*+}, and its two rules.
%   148  the functoriality square (34) for T.
%   149  the four-item recapitulation: Chern rings A^n, the A(u), the bilinear
%        A^p × A^q → A^{p+q}, and T : A^• → Z × (1 + completion).
%   150  A(T) = Z(T^) for a torus: C^i(ρ) the elementary symmetric functions of
%        the characters, C(ρ) = C(ρ')C(ρ''), the homomorphism into
%        Z × (1 + Z[[T^]]^+) injective with dense image, injectivity through
%        χ ↦ e^χ, density reduced to a statement about symmetric functions.
%   151  the Lemme that finishes it, then a pencil schema: Z[T^] onto
%        H^•(B_T,Z), A(T) below, and « Il faudrait prouver que Z[T^] → l'anneau
%        de Chern attaché à A(T) ≃ Z(T^) ».
%   152  K(G) for an algebraic group, CK(G) its ring of Chern classes, the
%        canonical c_G, the invariance under inner automorphisms, and the
%        characteristic homomorphism CK(G) → CK(X) = A(X) of a principal bundle.
%   153  the two compatibility diagrams (3) and (4); the Proposition computing
%        K(T) and CK(T) for a torus of dimension r.
%   154  the Weyl group: CK(G) → Z[T^]^W, not in general bijective, and (5 bis),
%        its k-linear analogue, bijective when the characteristic of k is prime
%        to the torsion of G. The run breaks off there.
%   156  C an abelian class of k-modules with exact forgetful functor and a set
%        of isomorphism representatives A_i; Lemme I identifying the
%        endomorphisms of the forgetful functor with the coherent families
%        (u_i), hence the complete linearly topologised ring U.
%   157  the canonical functor I into the « permitted » U-modules, and what is
%        wanted of it. The lower two thirds of the leaf are cancelled.
%   158  the two conditions, 1°) bijectivity and 2°) that k-homomorphisms be
%        exactly U-homomorphisms; the augmentation; a boxed statement on the
%        permitted U-modules. Heavily struck, with a cancelled block filling the
%        left margin.
%   159  two such algebras U, V: continuous homomorphisms V → U against
%        permitted functors C^U → C^V over C^k, with the underlined conclusion
%        that the two identify.
%   160  augmentations and the trivial module; then duality — a functor F^∨
%        putting a permitted U-module structure on the dual, with F^∨F^∨ = id
%        and the opposite ring U°. The batch ends mid-argument.
%
% Two material facts are stated once here rather than page by page. First, his
% hand in this run is a fast draft: the formulas carry and a good deal of the
% connective prose does not, so the apparatus is dense and \ill{} stands
% wherever the leaf will not support a reading. Second, several notes run
% slantwise up the left edge of pages 142, 144, 146 and 156 in a hand faster
% still; they are given as far as they go and no further.
%
% What the leaves do not support is marked and not filled: the parenthesis
% closing the first paragraph of page 142, the openings of the marginal notes,
% the cancelled block of page 157, the left-margin block of page 158, and the
% struck passages of pages 159 and 160.

\documentclass[11pt,a4paper]{article}
% Le préambule commun à toutes les transcriptions du fonds Grothendieck.
%
% Il tient en une idée : **l'incertitude se compose, elle ne se résout pas.**
% Une transcription qui lisse un mot douteux a détruit la seule chose pour
% laquelle on la faisait — un lecteur doit pouvoir savoir, à chaque endroit,
% ce qui était sur le feuillet et ce qui a été ajouté. Les six macros qui
% suivent sont donc l'appareil critique, et non de la décoration.
%
% Les mêmes commandes sont reconnues par `scripts/render.mjs`, qui produit la
% vue de lecture du site. Une macro ajoutée ici doit l'être aussi là-bas,
% faute de quoi le rendu échouera bruyamment — ce qui est voulu : un rendu
% silencieusement faux est pire qu'un rendu absent.

\ProvidesPackage{grothendieck}[2026/08/08 Transcriptions du fonds Grothendieck]

\usepackage{amsmath,amssymb,amsthm}
\usepackage{mathrsfs}
\usepackage{stmaryrd}
% Les diagrammes commutatifs sont partout dans ce fonds : c'est la langue
% même de la géométrie algébrique de Grothendieck, et une transcription qui
% les rendrait en prose ne transcrirait rien.
\usepackage{tikz-cd}
% Français par défaut — c'est la langue des feuillets — mais l'anglais est
% chargé aussi : la traduction et le résumé sont des documents anglais, et la
% typographie française (espaces avant les deux-points) y serait une faute.
% Ces documents basculent par \selectlanguage{english} en tête de corps.
\usepackage[english,french]{babel}
\usepackage{fontspec}
\usepackage{geometry}
\geometry{a4paper,margin=25mm,marginparwidth=28mm}
\usepackage{marginnote}
\usepackage{xcolor}
% ulem, not soul. `soul`'s \st and \ul are built on a character-by-character
% scan that XeLaTeX's font machinery breaks: the document fails with an
% undefined control sequence rather than degrading. ulem does the same two jobs
% and is Unicode-engine safe. `normalem` stops it hijacking \emph, which the
% transcriptions use for Grothendieck's own emphasis.
\usepackage[normalem]{ulem}
\usepackage{hyperref}
\hypersetup{colorlinks=true,linkcolor=black,urlcolor=black,pdfborder={0 0 0}}

% --- Métadonnées du lot -----------------------------------------------------
% Reprises telles quelles par le script de rendu, qui les affiche en tête de
% la vue de lecture. Elles fixent ce que le fichier prétend couvrir : sans
% elles, une transcription flotte sans point d'attache dans un fonds de seize
% mille feuillets.

\newcommand{\folder}[1]{\gdef\grothFolder{#1}}
\newcommand{\batch}[1]{\gdef\grothBatch{#1}}
\newcommand{\leaves}[2]{\gdef\grothFirst{#1}\gdef\grothLast{#2}}
\newcommand{\foldertitle}[1]{\gdef\grothFoldertitle{#1}}
\gdef\grothFolder{?}\gdef\grothBatch{?}\gdef\grothFirst{?}\gdef\grothLast{?}\gdef\grothFoldertitle{}

% --- Le résumé --------------------------------------------------------------
%
% Il ouvre la lecture modernisée, sous le titre « Résumé », et il est composé
% en petit. Ce n'est pas un ornement : c'est le seul passage du document qui
% ne soit pas des mathématiques, et un lecteur doit pouvoir le reconnaître
% avant de l'avoir lu — puis le sauter s'il connaît déjà le sujet, ou s'y
% arrêter s'il ne le connaît pas. Le filet qui le suit marque la reprise du
% corps.

\newenvironment{resume}
  {\small\setlength{\parskip}{0.45em}}
  {\par\medskip\hrule\medskip}

% --- Mots-clés ---------------------------------------------------------------
%
% En anglais, en clôture du résumé de la lecture modernisée : le vocabulaire
% moderne sous lequel chercher ce que les feuillets construisent. Le manifeste
% du site (`npm run manifest`) les extrait pour étiqueter la cote entière —
% cette ligne est la source unique des tags, il n'y a pas de fichier à côté.
\newcommand{\keywords}[1]{\par\smallskip\noindent
  {\footnotesize\textsc{Keywords} --- \textit{#1}}\par}

% --- L'appareil critique ----------------------------------------------------

% Le numéro de feuillet, en marge. C'est l'ancre qui permet de régler un
% désaccord sur une formule en regardant *un* feuillet plutôt qu'une chemise
% de six cents. Le site s'en sert aussi pour tourner les pages du fac-similé
% quand on fait défiler la transcription.
\newcommand{\leaf}[1]{%
  \marginnote{\footnotesize\color{gray!70}\textsf{#1}}%
  \label{leaf:#1}\ignorespaces}

% Illisible : on ne devine pas. Un mot inventé qui se lit comme les autres est
% le pire résultat possible.
\newcommand{\ill}{\textcolor{red!60!black}{[\,\ldots\,]}}

% Lecture douteuse : le mot est donné, et signalé comme incertain.
\newcommand{\uncertain}[1]{\uline{#1}}

% Ajout de l'éditeur — une lettre manquante, un mot sous-entendu.
\newcommand{\add}[1]{\textcolor{blue!50!black}{[#1]}}

% Biffé par Grothendieck lui-même. Ce qu'il a rayé fait partie du document :
% une rature dit souvent où la pensée a changé de direction.
\newcommand{\struck}[1]{\sout{#1}}

% Note du transcripteur, sur l'état matériel du feuillet.
\newcommand{\note}[1]{\par\smallskip{\small\color{gray!60!black}\textit{[#1]}}\par\smallskip}

% Note marginale de Grothendieck, distincte du corps du texte.
\newcommand{\marginal}[1]{\par\smallskip\noindent\hspace{1em}%
  {\small\color{gray!60!black}#1}\par\smallskip}

% --- Titre ------------------------------------------------------------------

\AtBeginDocument{%
  \begin{center}
    {\large\bfseries Fonds Alexandre Grothendieck --- cote n\textsuperscript{o}~\grothFolder}\\[2pt]
    {\itshape\grothFoldertitle}\\[4pt]
    {\small feuillets \grothFirst--\grothLast\ (lot \grothBatch)}\\[2pt]
    {\footnotesize\color{gray!60!black}Transcription établie d'après le fac-similé numérisé
     par l'Université de Montpellier.\\
     Les lectures incertaines sont soulignées, les illisibles notées [\,\ldots\,],
     les ajouts entre crochets.}
  \end{center}
  \vspace{1em}}

\endinput


\folder{10}
\batch{8}
\pages{141}{160}
\dating{[à partir de 1958]}
\watermark{Édition de démonstration}
\foldertitle{Catégories tensorielles : notes manuscrites (s.d.), tapuscrits (s.d.)}

\begin{document}

\section*{Anneaux des représentations de certains groupes algébriques (pages 142 à 148)}

\note{Le faisceau annoncé par la chemise de la page 140 commence ici. La page 141
est blanche. Les sept feuillets qui suivent sont de sa main, d'un trait rapide,
et il les numérote lui-même de (1) à (34).}

\page{142}

Considérons la catégorie des espaces vectoriels \add{de dim.\ finie} sur un corps
$k$ (qu'on fixe une fois pour toutes). Pour tout groupe fini $G$, soit $A^{G}$ la
catégorie des espaces vectoriels de dim.\ finie où $G$ opère. C'est aussi une
catégorie avec produit tensoriel \add{et involution} (et $A^{e}$ \ill{} est la
catégorie des \ill{}).

\struck{On a de plus un bifoncteur $A^{G} \times A^{G'} \to A^{G}$}

Cette catégorie dépend fonctoriellement de $G$ : si \ill{} un hom.\
$u : G \to G'$, \ill{} a un foncteur contravariant
\[ (1) \qquad A^{u} : A^{G'} \longrightarrow A^{G} \]
et
\[ (1') \qquad A^{\mathrm{id}_{G}} = 1_{A^{G}} , \qquad A^{uv} = A^{v} A^{u} . \]
\struck{Ces foncteurs sont \ill{}} Ils permettent de définir de façon \ill{} le
produit tensoriel extérieur
\[ (2) \qquad A^{G} \times A^{G'} \longrightarrow A^{G \times G'} \]
en posant, si $p_{1}, p_{2}$ \ill{} les projections de $G \times G'$ sur ses deux
facteurs,
\[ (2') \qquad M \boxtimes M' = \struck{\ill{}} \ A^{p_{1}}(M) \otimes A^{p_{2}}(M') . \]

\marginal{\ill{} dans \ill{} imparfait}

\ill{} la somme directe extérieure $\boxplus$ \ill{} « commutatif »,
« associatif ». $\boxtimes$ est un bifoncteur additif. \ill{} les structures
multiplicatives dans les $A^{G}$ \ill{} se reconstruire à partir de ces
dernières, par \ill{} la considération de $\Delta_{G} : G \to G \times G$ :
\[ (3) \qquad M \otimes M' = A^{\Delta_{G}}(M \boxtimes M') . \]
Enfin, on a aussi une notion de représentation induite, qui pour \ill{}
$u : G \to G'$ définit un foncteur covariant additif

\page{143}

\[ (4) \qquad I^{u} : A^{G} \longrightarrow A^{G'} \]
et
\[ (4') \qquad I^{\mathrm{id}_{G}} = 1_{A^{G}} , \qquad I^{uv} = I^{u} I^{v} . \]
\ill{} les compatibilités entre ces deux sortes de foncteurs \struck{\ill{}} et
les opérations $\otimes$ étant
\[ (5) \qquad A^{u}\bigl(I^{u}(M)\bigr) = M^{[G' : G]} \qquad
   (M \in A^{G} , \ G \subset G') \]
\struck{\ill{}}

On a d'ailleurs deux homomorphismes fonctoriels
\[ (6) \qquad M \longrightarrow A^{u} I^{u}(M) \]
\[ (7) \qquad M' \longrightarrow I^{u}\bigl(A^{u}(M')\bigr) \]
\ill{}

\marginal{(10) Si $g \in G$, et si on désigne par $\sigma_{g}$ l'automorphisme
intérieur \ill{}, on a pour $M \in A^{G}$ : $M \simeq I^{\sigma_{g}}(M) \simeq
A^{\sigma_{g}}(M)$, l'isomorphisme étant canonique et défini par \ill{}}

\[ (8) \qquad \text{si $u$ est un isom., alors } I^{u} = A^{u^{-1}} , \]
donc $I^{u}$ et $A^{u}$ sont des foncteurs inverses l'un de l'autre.

\struck{Définition} \ (9) \ Soit $s$ la symétrie $G_{1} \times G_{2} \to
G_{2} \times G_{1}$. Alors $I^{s}(M \boxtimes M') = M' \boxtimes M$.

Soit, pour un entier $p \geq 1$ et des entiers $p_{1}, \ldots, p_{h}$ de somme
$p$,
\[ (10) \qquad u_{p_{1}, \ldots, p_{h}} : \mathfrak{S}_{p_{1}} \times \cdots
   \times \mathfrak{S}_{p_{h}} \longrightarrow \mathfrak{S}_{p} \]
l'hom.\ (injectif) bien connu. Posons, pour $M \in A^{\mathfrak{S}_{p}}$,
$M' \in A^{\mathfrak{S}_{q}}$,
\[ (11) \qquad M * M' = I^{u_{p,q}}(M \boxtimes M') . \]
On a un isomorphisme canonique
\[ (12) \qquad M * M' \simeq M' * M , \]
\struck{\ill{}} \ill{} $\alpha$ \uncertain{unique} de $\mathfrak{S}_{p+q}$
vérifiant $\alpha\, u_{p,q}\, \alpha^{-1} = u_{q,p}\, s_{p,q}$, où $s_{p,q}$ est
la symétrie $\mathfrak{S}_{p} \times \mathfrak{S}_{q} \to \mathfrak{S}_{q} \times
\mathfrak{S}_{p}$, et un isomorphisme canonique
\[ (13) \qquad (M * M') * M'' \simeq M * (M' * M'') \simeq
   I^{u_{p,q,r}}(M \boxtimes M' \boxtimes M'') . \]

\page{144}

\struck{\ill{}}

Soit $M \in A^{G}$ ; alors pour tout entier $n$, $\otimes^{n} M$ est un objet de
$A^{G \times \mathfrak{S}_{n}} = (A^{G})^{\mathfrak{S}_{n}}$, qu'on note
$T^{n}(M)$. Ainsi $T^{n}$ est un foncteur (non plus additif !)
\[ (15) \qquad T^{n} : A^{G} \longrightarrow (A^{G})^{\mathfrak{S}_{n}} \]
dont la propriété essentielle est
\[ (16) \qquad \boxed{\ T^{n}(M + M') \simeq \sum_{p+q=n} T^{p}(M) * T^{q}(M')\ } \]

\marginal{expliquer 2 opérations : $\boxtimes M \in A^{\uncertain{G^{n}}}$ où on
considère $G^{n}$ et $\mathfrak{S}_{n}$ opère dans $\boxtimes M$ \ill{}}

De plus, le diagramme
\[ (17) \]

\begin{tikzcd}[column sep=large]
A^{G'} \arrow[r, "T^{n}_{G'}"] \arrow[d, "A^{u}"'] &
  (A^{G'})^{\mathfrak{S}_{n}} \arrow[d, "A^{u \times 1_{\mathfrak{S}_{n}}}"] \\
A^{G} \arrow[r, "T^{n}_{G}"'] & (A^{G})^{\mathfrak{S}_{n}}
\end{tikzcd}

est commutatif (à préciser les compatibilités \ill{} hom.\ fonctoriels).

Considérons enfin l'hom.
\[ (18) \qquad u_{p,q} : \mathfrak{S}_{p} \times \mathfrak{S}_{q}
   \longrightarrow \mathfrak{S}_{pq} \]
défini par
$u_{p,q}(\sigma, \tau) = \alpha_{p,q}(\sigma)\,
u_{\underbrace{\tau, \ldots, \tau}_{p}}\, \alpha_{p,q}(\sigma)^{-1}$,
où $\alpha_{p,q}(\sigma) \in \mathfrak{S}_{pq}$ est défini par
\[ (19) \qquad \begin{cases}
   \alpha_{p,q}(\sigma) \text{ est croissant dans les intervalles }
     \bigl] iq, (i+1)q \bigr] \\
   \text{et applique ceux-ci sur les intervalles }
     \bigl] \sigma(i) q, (\sigma(i)+1) q \bigr]
\end{cases} \]
Sous ces conditions, on a
\[ (20) \qquad T^{m}\bigl(T^{n}(M)\bigr) \simeq \struck{\ill{}}\
   A^{u_{m,n}}\bigl(T^{mn}(M)\bigr) \]
(les deux membres dans $A^{G \times \mathfrak{S}_{m} \times \mathfrak{S}_{n}}$).

\page{145}

On a aussi
\[ (21) \qquad T^{n}(M \otimes M') \simeq T^{n}(M) \otimes T^{n}(M')
   \ \in \ A^{G \times \mathfrak{S}_{n}} . \]

\note{Le feuillet ne porte que ces deux lignes.}

\page{146}

De plus, le foncteur $M * M'$ est bi-additif en chaque argument. Il y a un
élément unité pour $*$ :
\[ \mathbf{1} = K \in A^{\mathfrak{S}_{0}} \qquad
   \bigl( \text{par définition } \mathfrak{S}_{0} = \{e\} , \text{ et }
   u_{0,p} = u_{p,0} = \ill{} \bigr) . \]

Passons aux groupes de classes (pour fixer les idées, posons)
\[ (22) \qquad A(G) = K(A^{G}) . \]
$A(G)$ est un \struck{\ill{}} \emph{anneau de Chern} commutatif augmenté,
\ill{} ; $A(G)$ est un foncteur \struck{covariant} contravariant en $G$ :
$u : G \to G'$ définit
\[ (23) \qquad A(u) : A(G') \longrightarrow A(G) \]
et $A(1_{G}) = 1_{A(G)}$, $A(uv) = A(v) A(u)$. On définit, pour $\alpha \in A(G)$,
$\beta \in A(G')$,
\[ (24) \qquad \alpha \boxdot \beta = \struck{\ill{}}\
   A(p_{1})(\alpha) \cdot A(p_{2})(\beta) \ \in \ A(G \times G') . \]
D'ailleurs, \ill{} de deux structures de groupe sur les $A(G)$, et de l'opération
$\alpha \boxdot \beta$, c.-à-d.\ les structures d'anneau de $A(G)$ \ill{}

Il y a aussi un foncteur \add{additif}, pour $u : G \to G'$,
\[ (25) \qquad I(u) : A(G) \longrightarrow A(G') \]
(\ill{} mais pas multiplicatif). On a
\[ (26) \qquad \begin{cases}
   A(u) I(u) = \text{multiplication par } [G' : G] \\
   I(u) A(u) = \text{multiplication par } [G' : G]
\end{cases} \]

\note{Les deux lignes de (26) portent le même facteur $[G':G]$ sur le feuillet ;
la transcription les donne telles quelles.}

\page{147}

\ill{} avec $M * M'$ \ill{} Ces conditions \ill{} que si \ill{} est un
automorphisme intérieur, $A(u)$ et $I(u)$ sont \ill{} l'un de l'autre. Enfin
\[ (27) \qquad A(\sigma_{g}) = I(\sigma_{g}) = \text{identité de } A(G)
   \qquad (g \in G) . \]
Les propriétés soulignées peuvent \ill{} sur deux sommes directes :
\[ (28) \qquad A^{n} = A(\mathfrak{S}_{n}) \quad (n \geq 1) , \quad
   A^{0} = \mathbb{Z} \ ; \qquad A^{\bullet} = \coprod_{n=0}^{\infty} A^{n} \]
muni d'une structure d'anneau gradué commutatif. \ill{} (et on a de même les
$A(G)^{n} = A^{G \times \mathfrak{S}_{n}}$, $A(G)^{\bullet} =
\coprod_{0}^{\infty} A(G)^{n}$).

D'ailleurs, on a, pour $\alpha \in A^{p}$, $\beta \in A^{q}$,
\[ (29) \qquad A(u_{p,q})(\alpha * \beta) = \frac{(p+q)!}{p!\, q!}\
   \alpha \boxdot \beta \ \in \ A^{\mathfrak{S}_{p} \times \mathfrak{S}_{q}} , \]
d'où, si on \ill{} une multiplication dans $A^{\bullet}$ \ill{}
\[ (30) \qquad A(u_{p,p})(\alpha * \beta) = \frac{(2p)!}{p!\, q!}\
   \alpha \beta \]

\note{Dans (30), « $2p$ » est écrit au-dessus de « $p+q$ » rayé, le dénominateur
restant $p!\,q!$.}

qui relie la multiplication dans les $A^{n}$ \struck{\ill{}} à la multiplication
« globale » $*$.

De plus, la formule (16) \ill{} une forme \ill{} commode :
\[ (31) \qquad \begin{cases}
   T = T_{G} : A(G) \longrightarrow 1 + \widehat{A(G)^{\bullet}}^{\,+} \\
   T(\alpha) = \sum_{n} T^{n}(\alpha) \qquad (T^{0}(\alpha) = 1)
\end{cases} \]
définie par
\[ (32) \qquad T\bigl(\gamma(M)\bigr) = \sum_{n} T^{n}(M) \]
\[ (33) \qquad T^{n}(\alpha + \beta) = \sum_{p+q=n} T^{p}(\alpha) * T^{q}(\beta) \]
(on distingue \ill{} de (31) \ill{})
\[ (33) \qquad T^{n}(\alpha \beta) = T^{n}(\alpha)\, T^{n}(\beta) \]
(qui, comparé à (29), explique « pourquoi » \ill{})

\note{Deux formules portent ici le numéro (33).}

\page{148}

\ill{} $T^{n}(\alpha\beta)$ : il n'y a \ill{} de l'opération $*$ et des opérations
\ill{} à \ill{} multiplication. Il \ill{} sont commodes \ill{} les
$T^{n}(\lambda^{i}(\alpha))$, on y \uncertain{reviendra}. Bien entendu, les
homomorphismes $T_{G}$ sont \struck{\ill{}} fonctoriels : si $u : G \to G'$, le
diagramme
\[ (34) \]

\begin{tikzcd}[column sep=large]
A(G') \arrow[r, "T_{G'}"] \arrow[d, "A(u)"'] &
  1 + \widehat{A(G')^{\bullet}}^{\,+} \arrow[d, "\widehat{A(u)}"] \\
A(G) \arrow[r, "T_{G}"'] & 1 + \widehat{A(G)^{\bullet}}^{\,+}
\end{tikzcd}

\ill{}

\section*{Récapitulation (page 149)}

\note{Les quatre articles résument, au crayon, le formalisme des pages
précédentes. Le bas de la feuille ne porte rien.}

\page{149}

\begin{enumerate}
\item[(i)] Anneaux \struck{commutatifs} \add{de Chern \uncertain{involutifs}
      augmentés} $A^{n}$, $A^{0} = \mathbb{Z}$.

\item[(ii)] Pour tout hom.\ $u : \mathfrak{S}_{p} \to \mathfrak{S}_{q}$, un hom.\
      $A(u) : A^{q} \to A^{p}$, avec $A(uv) = A(v) A(u)$,
      $A(1_{\mathfrak{S}_{p}}) = 1_{A^{p}}$, et $A(\sigma_{g}) = 1_{A^{p}}$ si
      $g \in \mathfrak{S}_{p}$.

\item[(iii)] \struck{Des} applications bilinéaires
      $A^{p} \times A^{q} \to A^{p+q}$, $(\alpha, \beta) \mapsto \alpha * \beta$,
      faisant de $A^{\bullet} = \coprod_{0}^{\infty} A^{n}$ un anneau gradué
      commutatif avec unité. On a la relation
      \[ \frac{(2p)!}{(p!)^{2}}\, \alpha \beta = A(u_{p})(\alpha * \beta) \]
      où $u_{p}$ est le composé
      $\mathfrak{S}_{p} \xrightarrow{\ \Delta\ } \mathfrak{S}_{p} \times
      \mathfrak{S}_{p} \xrightarrow{\ u_{p,p}\ } \mathfrak{S}_{2p}$.

\item[(iv)] \struck{Des} homomorphismes de groupes
      \[ T : A^{\bullet} \longrightarrow \widetilde{A^{\bullet}}
         = \mathbb{Z} \times \bigl( 1 + \widehat{A^{\bullet +}} \bigr) \]
      défini par des composantes $T^{n}$, avec $T^{0}(\alpha) = 1$,
      $T^{1}(\alpha) = \alpha$, $T^{n}(\alpha) \in A^{n}$. On a donc
      \[ T^{n}(\alpha + \beta) = \sum_{p+q=n} T^{p}(\alpha)\, T^{q}(\alpha) \]
      \note{Le second facteur est écrit $T^{q}(\alpha)$ sur le feuillet, là où
      l'énoncé demande $T^{q}(\beta)$ ; la leçon est donnée telle quelle.}
      De plus
      \[ T^{n}(\alpha \cdot \beta) = T^{n}(\alpha)\, T^{n}(\beta) . \]
\end{enumerate}

\section*{$A(T) = \mathbb{Z}(\hat{T})$ (pages 150 et 151)}

\page{150}

Soit $\rho : T \to Gl(n)$ ; \ill{} induit $\rho : T \to k^{*n}$, d'où $n$
éléments $\alpha_{1}, \ldots, \alpha_{n}$ de $\hat{T}$. Les fonctions symétriques
élémentaires en les $\alpha_{i}$ sont notées $C^{i}(\rho)$
($0 \leq i \leq n$) ; elles ne dépendent que de $\rho$. Soit
$C(\rho) = \sum_{0}^{\infty} C^{i}(\rho) \in \mathbb{Z}[\hat{T}]$
(avec $C^{i}(\rho) = 0$ si $i > \deg \rho$) ; si $\rho$ est une extension de
$\rho'$ par $\rho''$, \ill{} $C(\rho) = C(\rho')\, C(\rho'')$.

On peut prolonger $\rho \mapsto (\deg \rho, C(\rho))$ en un hom.\ \ill{}
\[ (1) \qquad A(T) \simeq \mathbb{Z}(\hat{T}) \longrightarrow
   \mathbb{Z} \times \bigl( 1 + \mathbb{Z}[[\hat{T}]]^{+} \bigr) ; \]
c'est un hom.\ d'anneaux de Chern, caractérisé par ses valeurs sur les
$\rho_{\chi}$ ($\chi \in \hat{T}$) :
\[ \tilde{C}(\rho_{\chi}) = \bigl( 1 ,\ 1 + [\chi] \bigr) . \]
Je dis que cet hom.\ (1) est injectif \struck{\ill{}} et à image dense. Pour le
voir, on \ill{} avec \ill{}
\[ \varphi : \mathbb{Z} \times \bigl( 1 + \mathbb{Z}[[\hat{T}]]^{+} \bigr)
   \longrightarrow \mathbb{Q}[[\hat{T}]] \]
transformant $\chi$ en $e^{\chi}$, et on note que les $e^{\chi}$ sont
linéairement indépendants. Pour montrer que l'image est dense, on montre que si
\struck{\ill{}} des polynômes $C^{1}, \ldots, C^{n}$ sont donnés,
\struck{\ill{}} avec $C^{i} \in \mathbb{Z}[X_{1}, \ldots, X_{r}]$
($1 \leq i \leq n$), $C^{i}$ de degré $i$, on peut trouver des formes linéaires
$\alpha^{1}, \ldots, \alpha^{N}$ en les $X_{i}$, telles que
$C^{i} = \sigma_{i}(\alpha^{1}, \ldots, \alpha^{N})$ pour $1 \leq i \leq n$.
Récurrence sur

\page{151}

$n$, triviale pour $n = 0$. On \ill{} une récurrence \ill{} $C^{1} = \cdots =
C^{n-1} = 0$, donc à prouver le

\emph{Lemme.} Il existe des formes linéaires $\alpha^{1}, \ldots, \alpha^{N}$,
\ill{} des $X^{1}, \ldots, X^{r}$, telles que
\[ \sigma^{1}\bigl((\alpha^{i})\bigr) = \cdots =
   \sigma^{n-1}\bigl((\alpha^{i})\bigr) = 0 , \qquad
   \sigma^{n}\bigl((\alpha^{i})\bigr) = C^{n} \ \text{polynôme donné} . \]

\note{Le reste de la page est un schéma au crayon, dans la colonne de gauche,
séparé du corps par un trait vertical.}

$\mathbb{Z}[[\xi_{1}, \ldots]]$

\begin{tikzcd}
E_{T} \arrow[d] \\ B_{T}
\end{tikzcd}

\begin{tikzcd}[column sep=large]
\mathbb{Z}(\hat{T}) \arrow[r] & H^{*}(B_{T}, \mathbb{Z}) \\
A(T) \arrow[u, "\mathrm{alg}"] \arrow[ur, "\text{top.\ via Chern}"']
\end{tikzcd}

\note{Sur le feuillet, la flèche « alg » monte obliquement de $A(T)$ vers
$\mathbb{Z}(\hat{T})$ et la flèche « top.\ via Chern » monte verticalement de
$A(T)$ vers $H^{*}(B_{T},\mathbb{Z})$ ; le diagramme les redresse. La flèche
horizontale du haut est tracée à deux pointes.}

\begin{tikzcd}[column sep=large]
\mathbb{Z}[\hat{T}] \arrow[r, "\text{hom.\ car.}"] & H^{*}(X, \mathbb{Z}) \\
A(T) \arrow[u] \arrow[ur, "\text{top.\ via Chern}"']
\end{tikzcd}

Il faudrait prouver que $\mathbb{Z}[\hat{T}] \to$ l'anneau de Chern attaché à
$A(T)$ $\simeq \mathbb{Z}(\hat{T})$.

\section*{$K(G)$, $CK(G)$ et l'homomorphisme caractéristique (pages 152 à 154)}

\page{152}

Soit $G$ un groupe algébrique. On désignera par $K(G)$ l'anneau \struck{des}
engendré par les classes des représentations rationnelles irréductibles de $G$ :
en tant que groupe additif, il est de base libre ; c'est un \emph{anneau de
Chern}. On désignera par $CK(G)$ « l'anneau des classes de Chern » de $K(G)$
\add{c'est un anneau gradué}, \struck{\ill{}} et on a un homomorphisme canonique
\[ (1) \qquad K(G) \xrightarrow{\ c_{G}\ }
   \mathbb{Z} \times \bigl( 1 + CK(G)^{+} \bigr) . \]
$K(G)$ et $CK(G)$ sont des foncteurs en $G$, et (1) est un hom.\ fonctoriel. Si
$\sigma$ est un automorphisme intérieur de $G$, alors $\sigma^{*}$ est
l'automorphisme identique de $K(G)$ et $CK(G)$.

Soit $P$ un fibré algébrique principal de groupe structural $G$, de base $X$. À
toute représentation linéaire rationnelle $\rho$ de $G$ est associé un fibré
vectoriel \add{algébrique} sur $X$, d'où un élément de $K(X)$. On en conclut un
hom.\ d'anneaux de Chern
\[ K(G) \longrightarrow K(X) \]
d'où un hom.\ d'anneaux gradués
\[ (2) \qquad CK(G) \longrightarrow CK(X) = A(X) \]
appelé \emph{hom.\ caractéristique} du fibré $P$. Le \ill{}

\page{153}

\[ (3) \]

\begin{tikzcd}[column sep=large]
K(G) \arrow[r, "\text{ass.}"] \arrow[d, "c_{G}"'] & K(X) \arrow[d, "c_{X}"] \\
\widetilde{CK}(G) \arrow[r, "\text{hom.\ car.}"'] &
  \widetilde{CK}(X) = \widetilde{A}(X)
\end{tikzcd}

De plus, si \ill{} est un hom.\ rat.\ $G \xrightarrow{\ \varphi\ } G'$ et si $P'$
est associé à $P$ et $\varphi$, on a commutativité dans
\[ (4) \]

\begin{tikzcd}[column sep=large]
CK(G') \arrow[r, "\varphi^{*}"] \arrow[dr, "\tau_{P'}"'] & CK(G) \arrow[d, "\tau_{P}"] \\
 & A(X)
\end{tikzcd}

Il \uncertain{vaut} \ill{} calculer explicitement $CK(G)$.

Voici les résultats essentiels dans cette voie.

\emph{Proposition.} Soit $T$ un tore algébrique \add{de dim.\ $r$}. Alors
\[ K(T) \simeq \mathbb{Z}(\hat{T}) \qquad (\text{algèbre du groupe } \hat{T}) \]
\[ \bigl\{ \simeq \mathbb{Z}[x^{1}, \ldots, x^{r}]_{x^{1} \cdots x^{r}}
   \ \struck{\text{associé}}\ \add{\text{avec}}\ \varepsilon(x^{i}) = 1 ,
   \ \lambda x^{i} = 1 - x^{i} t \bigr\} \]
avec $\varepsilon(\chi) = 1$, $\lambda(\chi) = 1 + \chi t$ pour $\chi \in \hat{T}$.
On a $CK(T) \simeq \mathbb{Z}[\hat{T}]$ (algèbre symétrique sur $\mathbb{Z}$ du
$\mathbb{Z}$-module $\hat{T}$), avec
\[ c_{T}(\chi) = 1 + \chi \qquad \text{i.e.} \qquad
   c_{T}^{i}(\chi) = \begin{cases}
   1 & \text{si } i = 0 \\ \chi & \text{si } i = 1 \\ 0 & \text{si } i > 1
\end{cases} \]

\note{Le signe de $\lambda$ change d'une ligne à l'autre : $1 - x^{i}t$ dans
l'accolade, $1 + \chi t$ à la ligne suivante. Les deux leçons sont celles du
feuillet.}

\page{154}

Soit $G$ un groupe \struck{\ill{}} algébrique linéaire quelconque, $T$ un tore
maximal, $N$ son normalisateur et $W = N/T$. \struck{\ill{}} $N$ opère sur $G$ et
$T$ de façon compatible par automorphismes intérieurs, et \ill{} d'un \ill{}
canonique
\[ CK(G) \longrightarrow CK(T) \simeq \mathbb{Z}[\hat{T}] \]
est compatible avec les opérations de $N$. $N$ opère trivialement sur le premier
membre, et $T$ opère trivialement sur le second, donc on voit $W$ opère sur
$\mathbb{Z}[\hat{T}]$ via les opérations de $W$ sur $\hat{T}$. On a donc un hom.
\[ (5) \qquad CK(G) \longrightarrow \mathbb{Z}[\hat{T}]^{W} \]
où $\mathbb{Z}[\hat{T}]^{W}$ désigne \struck{le sous-groupe} \add{l'anneau} des
éléments de $\mathbb{Z}[\hat{T}]$ invariants sous $W$.

\struck{On} En général \struck{\ill{}} on n'a pas là une bijection. Mais c'est
\add{sans doute} une bijection s'il n'y a pas de torsion \struck{\ill{}}, et de
plus l'homomorphisme analogue
\[ (5\ \mathrm{bis}) \qquad C\bigl( K(G) \otimes_{\mathbb{Z}} k \bigr)
   \longrightarrow k[\hat{T}]^{W} \]
(où $k$ est un corps commutatif \ill{}) est bijectif si \struck{\ill{}}
\add{la car.\ de $k$ est} premier à la torsion de $G$.

\note{Le feuillet s'arrête là ; la page 155 est blanche.}

\section*{Classes abéliennes de $k$-modules et algèbres topologiques (pages 156 à 160)}

\note{Dernier faisceau du lot, de sa main et d'un trait plus rapide encore : les
formules portent, une bonne part de la prose de liaison ne porte pas, et les
quatre derniers feuillets sont largement raturés.}

\page{156}

$\mathcal{C}$ classe abélienne, \struck{\ill{}} dont les éléments sont des
\emph{modules} sur l'anneau commutatif $k$ et les homomorphismes des hom.\ de
modules, la composition des hom.\ étant la composition usuelle ; \ill{} que le
foncteur d'oubli $\mathcal{C} \to \mathcal{C}^{k}$ soit \emph{exact} et qu'il
existe un \emph{ensemble} de $A_{i} \in \mathcal{C}$ tels que \ill{}
$M \in \mathcal{C}$ soit isomorphe à l'un d'eux.

Notation $\operatorname{Hom}_{k}(M, N)$ si $M, N \in \mathcal{C}$.

Le noyau, l'image etc.\ \ill{} sont $k$ \ill{} : ce sont \ill{} ordinaires.

\marginal{\ill{}}

\emph{Lemme I.} \ $\mathcal{C} \to \mathcal{C}^{k}$ le foncteur \ill{}. Les
endomorphismes de ce foncteur d'oubli \add{\ill{}}
\struck{\ill{}} \ill{} les $k$-modules $(u_{i})$, $u_{i} \in L_{k}(A_{i})$, tels
que pour \ill{} $\varphi \in L_{k}(A_{i}, A_{j})$, on ait commutativité dans

\begin{tikzcd}
A_{i} \arrow[r, "\varphi"] \arrow[d, "u_{i}"'] & A_{j} \arrow[d, "u_{j}"] \\
A_{i} \arrow[r, "\varphi"'] & A_{j}
\end{tikzcd}

\[ \varphi u_{i} = u_{j} \varphi . \]

Ces endomorphismes forment donc un \ill{} \uncertain{sous-anneau} topologique de
$\prod_{i} L(A_{i})$, fermé pour la topologie produit, quand les $L(A_{i})$ sont
munis de la topologie \ill{} ($A_{i}$ étant discrets). Donc \ill{} muni d'une
topologie \struck{\ill{}} linéaire \ill{} \emph{complète}, qu'il \ill{}
$\mathcal{U}$. Si $V \in \mathcal{C}$, alors $V$ est un module sur $\mathcal{U}$,
et $\underline{\mathcal{U}} \to L_{s}(V)$ est continu ; de plus, les
$\varphi \in \operatorname{Hom}_{k}(V, W)$ sont des $\mathcal{U}$-homomorphismes.
\ill{}

\page{157}

\struck{\ill{}} On a ainsi un foncteur canonique, noté $I$, de $\mathcal{C}$ dans
les $\mathcal{U}$-modules « permis » (i.e.\ tels que
$\mathcal{U} \to L_{s}(V)$ soit continue). On se donne des conditions, pour que
ce soit \struck{bijectif} \add{un isomorphisme}. Si \ill{} les $A_{i}$ à
engendrement fini sur $k$, alors les $L(A_{i})$ sont \emph{discrets}, donc
$\mathcal{U}$ est un anneau \ill{} \emph{linéaire} \ill{} (avec une topologie
\ill{}) et \ill{} se borner à considérer les $\underline{\mathcal{U}}$-modules
permis à engendrement fini sur \struck{\ill{}} $k$. \ill{}

\struck{\ill{}}

\note{Les deux tiers inférieurs du feuillet sont annulés par de longues
diagonales et par une boucle tracée à la main ; le bloc ainsi rayé porte sur les
$\operatorname{Hom}_{k}(V,V)$, sur un foncteur $F_{\mathcal{U}}$ de
$\mathcal{C}^{\mathcal{U}}$ dans $\mathcal{C}$ et sur son identité avec $I$, et
ne rend ni ses phrases ni la plupart de ses formules.}

\page{158}

Pour tout $V \in \mathcal{C}^{k}$ ($k$-modules de type fini), soit $K(V)$ et
$K'(V)$ l'ensemble des « structures de type $\mathcal{C}$ » resp.\ des structures
de $\underline{\mathcal{U}}$-modules permis sur $k$. On a défini une \ill{}
fonctorielle $K(V) \longrightarrow K'(V)$. On veut :
\begin{enumerate}
\item[1°)] On veut que cet hom.\ soit bijectif ;
\item[2°)] Si $V, W \in \mathcal{C}$, on veut que les $k$-hom.\ de $V$ dans $W$
      soient identiques aux $\underline{\mathcal{U}}$-homomorphismes.
\end{enumerate}

On est amené à considérer les $\underline{\mathcal{U}}$-modules permis qui sont à
engendrement fini sur $k$. $\underline{\mathcal{U}}$ est un anneau
\struck{topologique} \add{linéairement \ill{} complet} \ill{} voisinages de
\ill{} \struck{\ill{}}

\ill{} anneaux de \ill{} type fini, s'obtenir ainsi.

Première structure supplémentaire : lorsque \ill{} $V$ libre \struck{\ill{}}
\ill{} qu'on doit être muni canoniquement d'une structure $k$, i.e.\ \ill{} sur
\ill{} une structure privilégiée de $\underline{\mathcal{U}}$-module, \ill{} il
\ill{} \ill{} sur $\underline{\mathcal{U}}$ une \emph{augmentation continue}
$\varepsilon$.

\[ \boxed{\ \text{Les } \underline{\mathcal{U}}\text{-modules permis sont les }
   \ill{} \text{ sur } A \ \ill{} \text{ qui sont de dim.\ finie sur } k . \ } \]

\note{Un bloc entier remplit la marge gauche du feuillet et est annulé par une
douzaine de diagonales ; il ne rend ni ses phrases ni ses formules.}

\page{159}

Soient $\underline{\mathcal{U}}, \underline{V}$ deux \struck{\ill{}}
« $\underline{k}$-algèbres », i.e.\ algèbres topologiques \struck{\ill{}}
satisfaisant aux conditions dites plus haut. Un hom.\ continu
\struck{compatible aux augmentations} $\underline{V} \to \underline{\mathcal{U}}$
définit un \struck{\ill{}} \add{foncteur} permis
$\mathcal{C}^{\underline{\mathcal{U}}} \to \mathcal{C}^{\underline{V}}$ ;
\struck{\ill{}} \ill{} tel que

\begin{tikzcd}[column sep=large]
\mathcal{C}^{\underline{\mathcal{U}}} \arrow[r]
  \arrow[dr, "I^{\underline{\mathcal{U}}}"'] &
  \mathcal{C}^{\underline{V}} \arrow[d, "I^{\underline{V}}"] \\
 & \mathcal{C}^{k}
\end{tikzcd}

soit commutatif.

Réciproquement, \add{\ill{}} \ill{} foncteur, défini \struck{\ill{}} soit
$V \in \mathcal{C}'{}^{\underline{\mathcal{U}}}$, \ill{} \ill{} $I(W)$,
\struck{\ill{}} \ill{} $F(V)$ \ill{} les $V$-modules. \ill{}

\note{Un bloc encadré à la main occupe le milieu du feuillet et est traversé par
deux longues diagonales ; il ne rend que les symboles isolés donnés ici.}

\ill{} $I^{\underline{\mathcal{U}}} = I^{\underline{V}} \circ F$, \ill{} de
$I^{\underline{V}}$ \ill{} de $I^{\underline{\mathcal{U}}}$, \ill{} de
$I^{\underline{V}}$ \ill{} $\underline{V}$ \ill{} $\underline{\mathcal{U}}$,
\ill{}

\emph{Ainsi, les hom.\ continus $\underline{V} \to \underline{\mathcal{U}}$
s'identifient aux hom.\ \uncertain{précédents} des foncteurs.}

\page{160}

Si \ill{} sur $\underline{\mathcal{U}}$ et $\underline{V}$ des augmentations
continues, les hom.\ compatibles avec l'augmentation correspondent aux foncteurs
permis, qui transforment le $\underline{\mathcal{U}}$-module trivial $k$ en le
$\underline{V}$-module trivial $k$. Transitivités, etc.

\struck{\ill{}}

\emph{Autres structures supplémentaires}

\struck{\ill{}} \ill{} donné un foncteur
$F^{\vee} : \mathcal{C}^{\underline{\mathcal{U}}} \to
\mathcal{C}^{\underline{\mathcal{U}}}$ \struck{\ill{}} associant à tout
$V \in \mathcal{C}^{\underline{\mathcal{U}}}$ une structure de
$\underline{\mathcal{U}}$-module permise sur le dual $V'$ de $V$,
\struck{De telle façon que} i.e.\ un foncteur tel que

\begin{tikzcd}[column sep=large]
\mathcal{C}^{\underline{\mathcal{U}}} \arrow[r, "F^{\vee}"] \arrow[d] &
  \mathcal{C}^{\underline{\mathcal{U}}} \arrow[d] \\
\mathcal{C}^{k} \arrow[r, "F^{\vee}"'] & \mathcal{C}^{k}
\end{tikzcd}

\[ \check{F}(V) = V' \]

soit commutatif. \struck{\ill{}} $V'$ \struck{\ill{} $\mathcal{U}^{\circ}$-module
(l'inverse de $\underline{\mathcal{U}}$)} \struck{\ill{}} \add{Ceci définit, pour}
\[ F^{\vee} F^{\vee} = \text{identité} . \]
\struck{\ill{}} soit alors $V$ \ill{} tout \ill{} permis sur
$\underline{\mathcal{U}}^{\circ}$ (l'inverse de $\underline{\mathcal{U}}$), $V'$
est un module sur $\underline{\mathcal{U}}$ (par \ill{}) donc $(V')'$ est un
module sur $\underline{\mathcal{U}}$. On a ainsi un \struck{\ill{}} foncteur
permis de $\underline{\mathcal{U}}$ \ill{}

\note{Le lot s'arrête là, en pleine phrase.}

\end{document}
